% Minimal retained preamble for lane-p2 cards taken from
% sources/logic-fl05-ramsey-like-theorems-for-the-schre/pinned-original-source.tex
% (arXiv:2412.11598v3, CC BY 4.0).  This source ships no separate preamble file:
% the pinned tex is a concatenation of the authors' main.tex and macros.tex and
% carries its preamble inline, so the blocks below are sliced mechanically from
% that file and nothing has been renamed or redefined.
%
% Copied verbatim from the source:
%   line 2        : the authors' class line
%   lines 3-5,7,9 : graphicx, amsthm, the ams bundle, cleveref, url
%   line 10       : tikz-cd, used by the authors' commutative diagrams
%   lines 20-108  : the authors' macro definitions, taken in full so that every
%                   macro of the retained spans is available
%   lines 113-124 : the authors' theorem declarations
%   lines 131-141 : the authors' machinery for restating a theorem
%   lines 143-148 : the authors' coloured comment macros
%
% Omitted source preamble lines and why:
%   line 1, 110-111 and 150 (concatenation comment, title, date, author): the
%       assembler supplies the title block, author list and date;
%   line 6 (hyperref) and line 8 (inputenc with the utf8x option): the forced
%       xelatex engine supplies the hyperref wiring itself and utf8x is
%       incompatible with it;
%   line 115 (numberwithin theorem by section): the assembler numbers the
%       retained material itself;
%   lines 120 and 121-123 are declarations that the authors left commented out or
%       that duplicate earlier ones and are not reproduced.
\documentclass{article}
\usepackage{graphicx}
\usepackage{amsthm}
\usepackage{amsfonts, amsmath, amssymb, amsthm}
\usepackage{cleveref}
\usepackage{url}
\usepackage{tikz-cd}
\newcommand{\Psf}{\mathsf{P}}
\newcommand{\Qsf}{\mathsf{Q}}
\newcommand{\Rsf}{\mathsf{R}}
\def\F{\mathcal{F}}
\def\U{\mathcal{U}}
\def\C{\mathcal{C}}
\def\I{\mathcal{I}}
\def\D{\mathcal{D}}
\def\L{\mathcal{L}}
\def\V{\mathcal{V}}
\def\A{\mathcal{A}}
\def\B{\mathcal{B}}
\def\P{\mathcal{P}}
\def\M{\mathcal{M}}
\def\Nc{\mathcal{N}}
\def\Q{\mathcal{Q}}
\newcommand{\cs}{2^{\NN}}
\newcommand{\bs}{\NN^{\NN}}
\newcommand{\bstr}{2^{<\NN}}
\newcommand{\uh}{\upharpoonright}
\newcommand{\converge}{\!\!\downarrow}
\newcommand{\diverge}{\!\!\uparrow}
\newcommand{\qvdash}{\operatorname{{?}{\vdash}}}
\newcommand{\nqvdash}{\operatorname{{?}{\nvdash}}}
\newcommand{\ISig}{\mathsf{I}\Sigma^0}
\newcommand{\Robinson}{\mathsf{Q}}
\newcommand{\IPi}{\mathsf{I}\Pi^0}
\newcommand{\IDelta}{\mathsf{I}\Delta^0}
\newcommand{\Iopen}{\mathsf{I_{open}}}
\newcommand{\BSig}{\mathsf{B}\Sigma^0}
\newcommand{\BPi}{\mathsf{B}\Pi^0}
\newcommand{\BDelta}{\mathsf{B}\Delta^0}
\newcommand{\BCSig}{\mathsf{BC}\Sigma^0}
\newcommand{\BCPi}{\mathsf{BC}\Pi^0}
\newcommand{\BCDelta}{\mathsf{BC}\Delta^0}
\newcommand{\LSig}{\mathsf{L}\Sigma^0}
\newcommand{\LPi}{\mathsf{L}\Pi^0}
\newcommand{\CSig}{\mathsf{C}\Sigma}
\newcommand{\RCA}[0]{\mathsf{RCA}}
\newcommand{\WKL}[0]{\mathsf{WKL}}
\newcommand{\WWKL}[0]{\mathsf{WWKL}}
\newcommand{\ACA}[0]{\mathsf{ACA}}
\newcommand{\ATR}[0]{\mathsf{ATR}}
\newcommand{\PIOOCA}[0]{\Pi^1_1\mathsf{-CA}}
\newcommand{\CAC}[0]{\mathsf{CAC}}
\newcommand{\ADS}[0]{\mathsf{ADS}}
\newcommand{\RRT}[0]{\mathsf{RRT}}
\newcommand{\FS}[0]{\mathsf{FS}}
\newcommand{\PFS}[0]{\mathsf{PFS}}
\newcommand{\TS}[0]{\mathsf{TS}}
\newcommand{\GRRT}[0]{\mathsf{GRRT}}
\newcommand{\LRRT}[0]{\mathsf{LRRT}}
\newcommand{\EM}[0]{\mathsf{EM}}
\newcommand{\SEM}[0]{\mathsf{SEM}}
\newcommand{\HalfEM}[0]{\mathsf{Half}\mbox{-}\mathsf{EM}}
\newcommand{\RT}[0]{\mathsf{RT}}
\newcommand{\SRT}[0]{\mathsf{SRT}}
\newcommand{\NN}[0]{\mathbb{N}}
\newcommand{\nat}[0]{\mathbb{N}}
\newcommand{\card}{\operatorname{card}}
\newcommand{\finsub}{\subseteq_{\mathtt{fin}}}
\newcommand{\W}{\mathrm{W}}
\newcommand{\sW}{\mathrm{sW}}
\newcommand{\lex}{\mathrm{lex}}
\newcommand{\base}{\operatorname{base}}
\newcommand{\range}{\operatorname{ran}}
\def\qt#1{``#1''}%

\newtheorem*{statement}{Statement}
\newtheorem{theorem}{Theorem}
\numberwithin{theorem}{section}
\newtheorem{maintheorem}[theorem]{Main Theorem}
\newtheorem{lemma}[theorem]{Lemma}
\newtheorem{question}[theorem]{Question}
\newtheorem{proposition}[theorem]{Proposition}
%\newtheorem{statement}[theorem]{Statement}
\newtheorem{remark}[theorem]{Remark}
\newtheorem{definition}[theorem]{Definition}
\newtheorem{corollary}[theorem]{Corollary}
\newtheorem{example}[theorem]{Example}

% Allow to repeat theorems
\makeatletter
\newtheorem*{rep@theorem}{\rep@title}
\newcommand{\newreptheorem}[2]{%
\newenvironment{rep#1}[1]{%
 \def\rep@title{#2 \ref{##1}}%
 \begin{rep@theorem}}%
 {\end{rep@theorem}}}
\makeatother

\newreptheorem{theorem}{Theorem}

% Comments
\usepackage{xcolor}
\newcommand{\ludovic}[1]{\textcolor{blue}{#1}}
\newcommand{\quentin}[1]{\textcolor{green}{#1}}
\newcommand{\lorenzo}[1]{\textcolor{orange}{#1}}
\newcommand{\oriola}[1]{\textcolor{red}{#1}}
\crefname{originalproposition}{original proposition}{original propositions}
\Crefname{originalproposition}{Original proposition}{Original propositions}
\makeatletter
\newlabel{prop:fs-n-k-strong-cone-avoidance}{{6.8}{source}{}{sourcepointer-1}{}}
\newlabel{prop:fs-n-k-strong-cone-avoidance@cref}{{[originalproposition][0][]6.8}{[1][0][]source}{}{}{}}
\makeatother
\title{Property (a) of a condition can be obtained for free by restricting the reservoir}
\author{Lorenzo Carlucci, Oriola Gjetaj, Quentin Le Houérou, Ludovic Levy Patey}
\date{arxiv:2412.11598v3; CC BY 4.0}
\begin{document}
\maketitle
\noindent\textit{Editorial note.} The counted unit is the article's own Lemma lem:fs-preserving-a at source lines 1121-1125 together with its proof at source lines 1126-1133, a straight adjacent pairing in the pinned file. The pinned file is a concatenation of the authors' main and macro files; it is the article's fixed version arXiv:2412.11598v3 under CC BY 4.0.\par\medskip
\noindent\textit{Editorial note.} This source ships no separate preamble file, so the retained preamble is sliced mechanically from the pinned file by the lane: the class line at source line 2, the package lines 3-5 and 7-10, the authors' macro block at lines 12-108 and their theorem declarations at lines 113-124, the restatement machinery at lines 131-141 and the coloured comment macros at lines 143-148. The title, date and author lines are supplied by the assembler, and the utf8x input encoding and hyperref are dropped because the forced xelatex engine supplies those itself.\par\medskip
\noindent\textit{Editorial note.} The retained statement and proof use the section's standing conventions, which the two context spans make available: the notion of a condition with its defining properties (a), (b) and (c) at source lines 1105-1117, and the convention at line 1119 that fixes the bounded-constrained coloring and its bounding function for the rest of the section. Both are context, not counted units, and neither overlaps the counted interval.\par\medskip
\noindent\textit{Editorial note.} The retained proof names one proposition of the same article, the strong cone-avoidance statement at source line 993, which is declared as a non-load-bearing navigation pointer. The article's labels are not part of any tag index, so the pointer's number was computed from the article's own section-wise numbering: it is 6.8, the eighth numbered environment of section 6, and the editorial record names the environments counted. Its source span points at the statement's own environment.\par\medskip
\noindent\textit{Editorial note.} The retained preamble keeps cleveref because the source writes some of its cross-references with the capitalised form, and keeps tikz-cd because other parts of the article draw diagrams; the retained span itself uses neither. No citation occurs in the retained spans, so no bibliography entry is supplied.\par\medskip
\noindent\textit{Allowed standard prerequisites (editorial).}\begin{itemize}
\item Conditions, bounded-constrained colorings and strong cone avoidance. The section works with conditions, each a triple of a coloring that is constrained by a computable bounding function, a finite stem and an infinite reservoir computing a fixed set D, together with an extension relation on them; the section also proves separately a strong cone-avoidance statement for the finite-dimensional free-set principle stating that for every finite tuple of colorings there is an infinite subset of a given reservoir that avoids the cone of D and is simultaneously free for all of them. The author's own definitions and that proposition are the allowed inputs.. The counted lemma is the step that shows the first defining property of a condition can always be recovered by shrinking the reservoir, which is what lets the rest of the diagonalisation work with arbitrary conditions.
\end{itemize}\medskip
\ifdefined\hypertarget\hypertarget{sourcepointer-1}{}\fi
\noindent\textit{Source reference.} Original proposition 6.8: sources/logic-fl05-ramsey-like-theorems-for-the-schre/pinned-original-source.tex, source lines 993--995. Navigation pointer invoked by the retained proof for the strong cone-avoidance statement that it iterates finitely many times; the cited statement is proved elsewhere in the same article and its proof is not reproduced inside the span. The article numbers theorem-like environments within the section, and counting the numbered environments of section 6 (Large Rainbow Ramsey Theorem avoids cones) gives this proposition the number 6.8, with the counted environments of that section at source lines 939, 945, 951, 964, 970, 978, 985, 993 and 1006; the number recorded here is that count, not a value copied from elsewhere.\par\medskip
\begin{definition}
A \emph{condition} is a triple $(f, \sigma, X)$ such that
$f \in \P$ is $b_f$-constrained, for some computable function $b_f : \NN \to \NN$,
$\sigma \in [\NN]^{<\omega}$, $X \subseteq \NN$ is an infinite set such that $\max \sigma < \min X$, and
\begin{enumerate}
    \item[(a)] for every~$s \in [\sigma \cup X]^{\leq h(\cdot)}$ with $\min s \in \sigma$, $f(s) \cap X \subseteq s$.
    \item[(b)] for every~$s \in [\sigma \cup X]^{\leq h(\cdot)}$, $f(s) \cap \sigma \subseteq s$.
    \item[(c)] $D \not \leq_T X$.
\end{enumerate}
A condition $d = (g, \tau, Y)$ \emph{extends} $c = (f, \sigma, X)$ (written $d \leq c$) if
$g \leq f$, $\tau \succeq \sigma$, $Y \subseteq X$ and $\tau \setminus \sigma \subseteq X$.
\end{definition}


In what follows, fix a coloring $f \in \P$ which is $b_f$-constrained, for some computable function $b_f : \NN \to \NN$.

\begin{lemma}\label[lemma]{lem:fs-preserving-a}
For every $\sigma \in [\NN]^{<\omega}$ and every infinite set~$X$ such that $\max \sigma < \min X$ and $D \not \leq_T X$,
there is an infinite set $Y \subseteq X$ such that $D \not \leq_T Y$ and 
for every~$s \in [\sigma \cup Y]^{\leq h(\cdot)}$ with $\min s \in \sigma$, $f(s) \cap Y \subseteq s$.
\end{lemma}

\begin{proof}
For every~$t \in [\sigma]^{\leq h(\cdot)}$ with $t \neq \emptyset$, let $f_t : [X]^{\leq h(\min t) - |t|} \to [\NN]^{\leq b_f(\min t)}$ be defined for every~$u$ by $f_t(u) = f(t \cdot u)$.
By finitely many successive applications of \Cref{prop:fs-n-k-strong-cone-avoidance},
there exists an infinite subset~$Y \subseteq X$ such that $D \not \leq_T Y$
and such that $Y$ is simultaneously $f_t$-free for every~$t \in [\sigma]^{\leq h(\cdot)}$ with $t \neq \emptyset$.

We claim that $Y$ is our desired set. Indeed, for every~$s \in [\sigma \cup Y]^{\leq h(\cdot)}$ with $\min s \in \sigma$, letting $t = s \cap \sigma$ and $u = s \cap Y$, we have $f(s) = f_t(u)$, so by $f_t$-freeness of~$Y$, $f(s) = f_t(u) \cap Y \subseteq u \subseteq s$.
\end{proof}

\end{document}
